\documentclass[a4paper,fontset=none]{ctexart}

\usepackage{anyfontsize}
\usepackage{amsmath,mathrsfs,wasysym}
\usepackage{graphicx,caption,subcaption}
\captionsetup{font=small}
\usepackage{geometry,parskip}
\geometry{top=3cm,bottom=3cm,left=2.75cm,right=2.75cm}
\usepackage{indentfirst}
\setlength{\parindent}{0em}
\usepackage[T1]{fontenc}
\usepackage[libertine,vvarbb]{newtxmath}
\usepackage[scr=rsfso,frak=euler,bb=ams]{mathalfa}
\usepackage{bm}
\usepackage{fontspec}
\setmainfont{Linux Libertine O}
\setsansfont{Linux Biolinum O}
\setCJKmainfont{SimSun}[BoldFont=Noto Sans CJK SC Medium]

\begin{document}

    \title{\textbf{实验二十六\ \ 真空镀膜}}
    \author{1900011604 张植竣}
    \date{2021年4月10日}

    \maketitle

    \section*{1\ \ 实验数据和现象}

    \subsection*{1.1\ \ 估计本实验条件下蒸镀薄膜所需真空度的下限}

    由气体分子的平均自由程：
    \[
    \overline{\lambda} = \frac{kT}{\sqrt{2}\pi d^2 p}
    \]
    得真空度的下限（即压强$p$的上限）：
    \[
    p_\mathrm{max} = \frac{kT}{\sqrt{2}\pi d^2 \overline{\lambda}}
    \]
    其中$T$和$p$是铜蒸发镀膜时的温度和压强，$d$是铜原子的直径，取$d=256\,\mathrm{pm}$。铜的沸点是$1084.62\,\text{\textcelsius}$，但由于压强降低，估计沸点为$1000\,\text{\textcelsius} = 1273.15\,\mathrm{K}$。气体分子的平均自由程$\overline{\lambda}$要远大于蒸发源到衬底的距离$x$，估计$x \approx 8\,\mathrm{cm}$，并取$\overline{\lambda} = 50x = 4\,\mathrm{m}$。这给出：
    \[
    p_\mathrm{max} = \frac{kT}{\sqrt{2}\pi d^2 \overline{\lambda}} = 15.09\,\mathrm{mPa}
    \]
    这与实验要求的$13\,\mathrm{mPa}$较为符合。

    \newpage
    \subsection*{1.2\ \ 分子泵开始工作后系统压强$p$随时间$t$的变化关系}

    以下是分子泵开始工作后系统压强$p$随时间$t$变化关系的数据表：（$t=0\,\mathrm{min}$时$p=85.2\mathrm{mPa}$）

    \begin{center}
    \begin{minipage}{0.45\textwidth}
        \begin{tabular}{|c|c|c|c|c|c|}
            \hline
            $t/\mathrm{min}$ & 0.5 & 1.0 & 1.5 & 2.0 & 2.5 \\
            \hline
            $p/\mathrm{mPa}$ & 76.6 & 70.0 & 64.3 & 59.6 & 55.4 \\
            \hline
            $t/\mathrm{min}$ & 3.0 & 3.5 & 4.0 & 4.5 & 5.0 \\
            \hline
            $p/\mathrm{mPa}$ & 51.4 & 48.5 & 45.6 & 43.1 & 40.8 \\
            \hline
            $t/\mathrm{min}$ & 5.5 & 6.0 & 6.5 & 7.0 & 7.5 \\
            \hline
            $p/\mathrm{mPa}$ & 38.6 & 36.8 & 35.1 & 33.6 & 32.2 \\
            \hline
            $t/\mathrm{min}$ & 8.0 & 8.5 & 9.0 & 9.5 & 10.0 \\
            \hline
            $p/\mathrm{mPa}$ & 30.9 & 29.7 & 28.5 & 27.6 & 26.7 \\
            \hline
            $t/\mathrm{min}$ & 10.5 & 11.0 & 11.5 & 12.0 & 12.5 \\
            \hline
            $p/\mathrm{mPa}$ & 26.0 & 25.3 & 24.7 & 23.9 & 23.3 \\
            \hline
            $t/\mathrm{min}$ & 13.0 & 13.5 & 14.0 & 14.5 & 15.0 \\
            \hline
            $p/\mathrm{mPa}$ & 22.6 & 22.2 & 21.7 & 21.1 & 20.7 \\
            \hline
            $t/\mathrm{min}$ & 15.5 & 16.0 & 16.5 & 17.0 & 17.5 \\
            \hline
            $p/\mathrm{mPa}$ & 20.2 & 19.7 & 19.4 & 19.1 & 18.7 \\
            \hline
            $t/\mathrm{min}$ & 18.0 & 18.5 & 19.0 & 19.5 & 20.0 \\
            \hline
            $p/\mathrm{mPa}$ & 18.3 & 17.9 & 17.6 & 17.3 & 16.9 \\
            \hline
            $t/\mathrm{min}$ & 20.5 & 21.0 & 21.5 & 22.0 & 22.5 \\
            \hline
            $p/\mathrm{mPa}$ & 16.7 & 16.4 & 16.1 & 15.8 & 15.6 \\
            \hline
            $t/\mathrm{min}$ & 23.0 & 23.5 & 24.0 & 24.5 & 25.0 \\
            \hline
            $p/\mathrm{mPa}$ & 15.4 & 15.2 & 14.9 & 14.7 & 14.5 \\
            \hline
        \end{tabular}
    \end{minipage}
    \begin{minipage}{0.45\textwidth}
        \begin{tabular}{|c|c|c|c|c|c|}
            \hline
            $t/\mathrm{min}$ & 25.5 & 26.0 & 26.5 & 27.0 & 27.5 \\
            \hline
            $p/\mathrm{mPa}$ & 14.3 & 14.1 & 13.9 & 13.7 & 13.6 \\
            \hline
            $t/\mathrm{min}$ & 28.0 & 28.5 & 29.0 & 29.5 & 30.0 \\
            \hline
            $p/\mathrm{mPa}$ & 13.5 & 13.3 & 13.2 & 13.1 & 13.0 \\
            \hline
            $t/\mathrm{min}$ & 30.5 & 31.0 & 31.5 & 32.0 & 32.5 \\
            \hline
            $p/\mathrm{mPa}$ & 12.8 & 12.7 & 12.6 & 12.5 & 12.4 \\
            \hline
            $t/\mathrm{min}$ & 33.0 & 33.5 & 34.0 & 34.5 & 35.0 \\
            \hline
            $p/\mathrm{mPa}$ & 12.3 & 12.2 & 12.1 & 11.9 & 11.8 \\
            \hline
            $t/\mathrm{min}$ & 35.5 & 36.0 & 36.5 & 37.0 & 37.5 \\
            \hline
            $p/\mathrm{mPa}$ & 11.7 & 11.6 & 11.5 & 11.4 & 11.3 \\
            \hline
            $t/\mathrm{min}$ & 38.0 & 38.5 & 39.0 & 39.5 & 40.0 \\
            \hline
            $p/\mathrm{mPa}$ & 11.2 & 11.1 & 11.0 & 10.9 & 10.8 \\
            \hline
            $t/\mathrm{min}$ & 40.5 & 41.0 & 41.5 & 42.0 & 42.5 \\
            \hline
            $p/\mathrm{mPa}$ & 10.7 & 10.6 & 10.5 & 10.4 & 10.3 \\
            \hline
            $t/\mathrm{min}$ & 43.0 & 43.5 & 44.0 & 44.5 & 45.0 \\
            \hline
            $p/\mathrm{mPa}$ & 10.3 & 10.2 & 10.2 & 10.1 & 10.0 \\
            \hline
            $t/\mathrm{min}$ & 45.5 & 46.0 & 46.5 & 47.0 & 47.5 \\
            \hline
            $p/\mathrm{mPa}$ & 9.9 & 9.9 & 9.8 & 9.7 & 9.6 \\
            \hline
            $t/\mathrm{min}$ & 48.0 & 48.5 & 49.0 & 49.5 & 50.0 \\
            \hline
            $p/\mathrm{mPa}$ & 9.6 & 9.5 & 9.4 & 9.4 & 9.3 \\
            \hline
        \end{tabular}
    \end{minipage}
    \end{center}

    作$p-t$图像如下：

    \begin{figure}[!ht]
        \centering
        \includegraphics[width=0.75\textwidth]{Figure_1.png}
        \caption{分子泵开始工作后系统压强$p(\mathrm{mPa})$随时间$t(\mathrm{min})$的变化关系}
    \end{figure}

    可见压强始终降低，且下降的速度先快后慢，最后趋向于一定值（极限压强）。用计算器拟合数据（不包括$t = 0\,\mathrm{min}$时的数据），最佳模型是$p = at^{-b}$。拟合结果为：
    \[
    a = 96.853,\qquad b = 0.587,\qquad |r| = 0.993
    \]
    而指数衰减模型$p \sim p_0 + C\mathrm{e}^{-\beta t}$和$p \sim \dfrac{1}{t}$的模型都不够理想。实际上压强的变化规律较为复杂，介于$p = p_0 + at^{-b}$和$p \sim p_0 + C\mathrm{e}^{-\beta t}$之间。

    \subsection*{1.3\ \ 加预蒸发电流，记录电流的变化情况}

    \textbf{预蒸发电流$I=5\,\mathrm{A}$}

    电流无明显变化。

    \textbf{预蒸发电流$I=15\,\mathrm{A}$}

    通电后不久电流先稍有减小，应是因为通电发热使电阻丝温度升高，电阻率增大；经过人工干预之后稳定在$I=15\,\mathrm{A}$。

    \subsection*{1.4\ \ 维持预蒸发电流大小不变，记录压强的变化情况}

    \textbf{预蒸发电流$I=5\,\mathrm{A}$，$t=0\,\mathrm{s}$时初始压强$p=8.8\,\mathrm{mPa}$}

    \begin{center}
    \begin{tabular}{|c|c|c|c|c|c|c|c|c|}
        \hline
        $t/\mathrm{s}$ & 3 & 6 & 9 & 12 & 15 & 18 & 21 & 24 \\
        \hline
        $p/\mathrm{mPa}$ & 8.8 & 8.9 & 8.9 & 9.0 & 9.0 & 9.1 & 9.1 & 9.1 \\
        \hline
        $t/\mathrm{s}$ & 27 & 30 & 33 & 36 & 39 & 42 & 45 & 48 \\
        \hline
        $p/\mathrm{mPa}$ & 9.1 & 9.1 & 9.2 & 9.2 & 9.3 & 9.4 & 9.4 & 9.5 \\
        \hline
        $t/\mathrm{s}$ & 51 & 54 & 57 & 60 & 63 & 66 & 69 & 72 \\
        \hline
        $p/\mathrm{mPa}$ & 9.6 & 9.7 & 9.8 & 9.9 & 10.1 & 10.2 & 10.3 & 10.2 \\
        \hline
        $t/\mathrm{s}$ & 75 & 78 & 81 & 84 & 87 & 90 & 93 & 96 \\
        \hline
        $p/\mathrm{mPa}$ & 10.2 & 10.2 & 10.2 & 10.0 & 10.0 & 9.9 & 9.8 & 9.8 \\
        \hline
        $t/\mathrm{s}$ & 99 & 102 & 105 & 108 & 111 & 114 & 117 & 120 \\
        \hline
        $p/\mathrm{mPa}$ & 9.7 & 9.7 & 9.7 & 9.6 & 9.6 & 9.6 & 9.6 & 9.6 \\
        \hline
    \end{tabular}
    \end{center}

    作图如下：

    \begin{figure}[!ht]
        \centering
        \includegraphics[width=0.75\textwidth]{Figure_3.png}
        \caption{预蒸发电流$I=5\,\mathrm{A}$时，系统压强$p(\mathrm{mPa})$随时间$t(\mathrm{s})$的变化关系}
    \end{figure}

    可见系统压强先增大后减小。开始时系统压强增大是因为电阻丝和铜丝通电发热，表面的杂质粒子蒸发为气态，蒸发速率大于分子泵吸气的速率；之后压强减小是因为杂质粒子大部分被分子泵抽走，蒸发速率小于分子泵吸气的速率。

    \textbf{预蒸发电流$I=15\,\mathrm{A}$，$t=0\,\mathrm{s}$时初始压强$p=8.6\,\mathrm{mPa}$}

    \begin{center}
    \begin{tabular}{|c|c|c|c|c|c|c|c|c|}
        \hline
        $t/\mathrm{s}$ & 3 & 6 & 9 & 12 & 15 & 18 & 21 & 24 \\
        \hline
        $p/\mathrm{mPa}$ & 9.3 & 12.5 & 14.9 & 14.4 & 19.2 & 35.4 & 15.2 & 14.3 \\
        \hline
        $t/\mathrm{s}$ & 27 & 30 & 33 & 36 & 39 & 42 & 45 & 48 \\
        \hline
        $p/\mathrm{mPa}$ & 14.7 & 15.8 & 17.0 & 18.8 & 20.4 & 22.2 & 23.8 & 25.3 \\
        \hline
        $t/\mathrm{s}$ & 51 & 54 & 57 & 60 & 63 & 66 & 69 & 72 \\
        \hline
        $p/\mathrm{mPa}$ & 26.6 & 27.7 & 28.8 & 29.5 & 30.1 & 30.5 & 30.8 & 30.9 \\
        \hline
        $t/\mathrm{s}$ & 75 & 78 & 81 & 84 & 87 & 90 & 93 & 96 \\
        \hline
        $p/\mathrm{mPa}$ & 31.0 & 31.0 & 30.8 & 30.6 & 30.4 & 30.1 & 29.8 & 29.5 \\
        \hline
        $t/\mathrm{s}$ & 99 & 102 & 105 & 108 & 111 & 114 & 117 & 120 \\
        \hline
        $p/\mathrm{mPa}$ & 29.2 & 29.0 & 28.7 & 28.5 & 28.2 & 27.9 & 27.6 & 27.4 \\
        \hline
    \end{tabular}
    \end{center}

    作图如下：

    \begin{figure}[!ht]
        \centering
        \includegraphics[width=0.75\textwidth]{Figure_4.png}
        \caption{预蒸发电流$I=15\,\mathrm{A}$时，系统压强$p(\mathrm{mPa})$随时间$t(\mathrm{s})$的变化关系}
    \end{figure}

    可见系统压强先突然增大再突然减小，之后先较缓慢增大后缓慢减小。开始时系统压强突变是因为电阻丝通电发热使得铜丝熔化，电路电阻突然减小，电流突然增大使得发热量突然增大，部分铜在此时蒸发为气态，系统压强增大，此时立即减小电流以降低铜的蒸发速率，系统压强也跟着减小；调整电流使其稳定在$I=15\,\mathrm{A}$，之后压强先增大也是因为铜丝受热蒸发，达到该电流下蒸发速率的最大值后，分子泵吸气速率大于蒸发速率使得压强减小。

    \subsection*{1.5\ \ 加蒸发电流，记录成膜过程和系统最大压强值}

    成膜过程：刚通电时成膜速度非常缓慢，镀膜不明显；之后铜突然大量蒸发，伴随耀眼的橙色闪光，约$2\sim3\,\mathrm{s}$即可看到玻片表面沉积一层深色带金属光泽的薄膜。

    系统最大压强值为$p = 73.8\,\mathrm{mPa}$，估算此时铜原子的平均自由程为：
    \[
    \overline{\lambda} = \frac{kT}{\sqrt{2}\pi d^2 p} = 0.818\,\mathrm{m}
    \]
    其余参数与\textbf{1.1}节中相同，即$T = 1273.15\,\mathrm{K}$，$d = 256\,\mathrm{pm}$。

    \subsection*{1.6\ \ 第二次镀膜}

    以下是分子泵开始工作后系统压强$p$随时间$t$变化关系的数据表：（$t=0\,\mathrm{min}$时$p=68.4\mathrm{mPa}$）

    \begin{center}
    \begin{tabular}{|c|c|c|c|c|c|c|}
        \hline
        $t/\mathrm{min}$ & 0.5 & 1.0 & 1.5 & 2.0 & 2.5 & 3.0 \\
        \hline
        $p/\mathrm{mPa}$ & 62.8 & 57.8 & 53.6 & 49.6 & 46.6 & 43.4 \\
        \hline
        $t/\mathrm{min}$ & 3.5 & 4.0 & 4.5 & 5.0 & 5.5 & 6.0 \\
        \hline
        $p/\mathrm{mPa}$ & 40.7 & 38.2 & 35.8 & 33.9 & 32.3 & 30.6 \\
        \hline
        $t/\mathrm{min}$ & 6.5 & 7.0 & 7.5 & 8.0 & 8.5 & 9.0 \\
        \hline
        $p/\mathrm{mPa}$ & 29.2 & 27.6 & 26.3 & 25.3 & 24.4 & 23.5 \\
        \hline
        $t/\mathrm{min}$ & 9.5 & 10.0 & 10.5 & 11.0 & 11.5 & 12.0 \\
        \hline
        $p/\mathrm{mPa}$ & 22.8 & 22.0 & 21.3 & 20.6 & 20.0 & 19.4 \\
        \hline
        $t/\mathrm{min}$ & 12.5 & 13.0 & 13.5 & 14.0 & 14.5 & 15.0 \\
        \hline
        $p/\mathrm{mPa}$ & 18.9 & 18.3 & 17.8 & 17.2 & 16.9 & 16.5 \\
        \hline
    \end{tabular}
    \end{center}

    \begin{figure}[!ht]
        \centering
        \includegraphics[width=0.75\textwidth]{Figure_2.png}
        \caption{第二次镀膜时分子泵开始工作后系统压强$p(\mathrm{mPa})$随时间$t(\mathrm{min})$的变化关系}
    \end{figure}

    可以看出第二次镀膜时分子泵抽真空过程中压强下降的速度略快于第一次镀膜时抽真空的速度（例如：第一次镀膜时压强从$68.4\,\mathrm{mPa}$下降到$16.5\,\mathrm{mPa}$用了大约$19.7\,\mathrm{min}$，第二次镀膜时只用了约$15.0\,\mathrm{min}$）。这是因为第一次抽真空时将分子泵上附着的气体分子也一并抽去，第二次抽真空时分子泵上附着的气体分子较少，抽气更快。

    \newpage
    \subsection*{1.7\ \ 镀膜成果}

    成功在两块玻片和一块1元硬币上镀膜，成果如下：

    \begin{figure}[!ht]
        \centering
        \begin{subfigure}[t]{0.4\textwidth}
            \centering
            \includegraphics[width=1\textwidth]{Photo_1.jpg}
            \subcaption{第一次镀膜}
        \end{subfigure}{\tiny}
        \quad
        \begin{subfigure}[t]{0.4\textwidth}
            \centering
            \includegraphics[width=1\textwidth]{Photo_2.jpg}
            \subcaption{第二次镀膜}
        \end{subfigure}
        \caption{镀膜成果}
    \end{figure}

    \newpage
    \section*{2\ \ 收获与感想}

    \subsection*{真空镀膜的经验}

    \begin{itemize}
        \item 保持玻片的光洁，虽然实验室没有提供去离子水、酒精和超声波清洗机，但是还是应该用洗耳球和擦镜纸清理干净表面的灰尘和污渍，并用镊子拿取玻片，镀膜前不得用手接触玻片。

        \item 第一次抽真空时将分子泵上附着的气体分子也一并抽去，第二次抽真空时分子泵上附着的气体分子较少。这也说明分子泵表面附着气体分子的速度较慢。两次镀膜的间隔时间越短，第二次抽真空时所用的时间就越短。

        \item 如果不用自动开启”高真空“档位的功能，而改用手动先开启”低真空“挡，抽至极限压强再切换”高真空“挡（前提是”低真空“档位的真空计显示正常），可以加快抽气速度。

        \item 镀上的铜膜极薄，容易磨损，取用时不能用手或者镊子，只能用纸片承载。
    \end{itemize}

\end{document}
